\section{Problema}
A metodologia expositiva tem origem no Mundo Antigo \cite{andreataAulaExpositivaPaulo2019}. Nesse sentido, é evidente que perseverou por incontáveis séculos, atravessando diferentes paradigmas tecnológicos e perdurando até o contexto atual. É uma técnica na qual o emissor é protagonista no processo de ensino, enquanto o ouvinte --- que devia receber maior atenção para a retenção do conteúdo --- é relegado a coadjuvante nesse vital processo.

Essa problemática é constantemente levantada por acadêmicos, como é o caso de \citeonline{vasconcellosMetodologiaDialeticaEm1992} e \citeonline{pereiraCriticaMetodologiaTradicional2022}, que apontam o mero papel de ouvinte como prejudicial aos estudantes, que não são estimulados a engajar; tem uma postura meramente passiva. Eles ainda citam que o docente procura verificar se todos estão entendendo, mas por meio de perguntas vagas, indagando os estudantes quanto à sua compreensão. Nesse processo, os estudantes apenas concordam, e a conversa é finalizada no ato. Não há acompanhamento ativo, apenas uma estrutura hierarquizada e de efetividade discutível.

Os estudos de \citeonline{vasconcellosMetodologiaDialeticaEm1992} e \citeonline{pereiraCriticaMetodologiaTradicional2022} foram feitos em ambiente presencial de ensino, no  qual há mais facilidade de interação entre as partes, além de melhor controle sobre a atenção e engajamento da turma. Então, se mesmo nessas condições ambos os estudos se mostram céticos quanto à efetividade da aula expositiva, é de se esperar que em ambiente remoto esse cenário seja ainda pior.

Com a pandemia de COVID-19, o ensino a distância ganhou ainda mais espaço, especialmente no Ensino Superior, que já apresentava tendência de crescimento exponencial no Brasil \cite{moranEnsinoSuperiorDistancia2009}. Apesar de ser um país relativamente pobre e com acesso limitado à tecnologia, o Brasil registra uma proporção significativa de estudantes universitários matriculados em cursos a distância, tanto totalmente online quanto semipresenciais.
